\clearpage
\subsection{等变情形}\label{sec:17k}
先作一些简单的几何准备。设 $M$ 为带 $\Z/2$ 作用 $\iota_M$ 的正合辛流形。
\begin{definition*}
$M$ 中的\emph{不变带标架正合 Lagrangian 球面}，是一个 $\iota_M$-不变的正合 Lagrangian 子流形 $V$，加上等变微分同胚 $v:S\to V$；这里 $S$ 是某个带线性 $\Z/2$ 作用的内积空间中的单位球面。
\end{definition*}
由于 Lagrangian 管状邻域定理可以与有限群作用相容，而 Dehn 扭转的局部模型关于 $O_{n+1}$ 等变，故沿这样的球面的 Dehn 扭转 $\tau_V$ 自然地成为一个等变正合辛自同构。给定带逐纤维对合的正合 Lefschetz 纤维化，其相应平行移动映射等变，消失圈为等变带标架球面，并且 Picard–Lefschetz 公式\eref{16.6}的适当等变版本成立。反过来，给定 $M$ 中一个或多个等变带标架球面，可以使\stref{引理}{16.8}和\stref{引理}{16.9}中的 Lefschetz 纤维化带有逐纤维对合。

现在假设处于\secref{14e}的情形，从而可定义等变 Fukaya 范畴。将其记为 $\A=\F(M)^{\mathrm{naive}}$，将其上的 $\Z/2$ 作用记为 $\eps_{\A}$。假设有一个等变带标架正合 Lagrangian
\SourcePage{272}
球面 $V=L_0$，并可将其赋予等变膜结构 $L_0^\#=(L_0,\alpha_0^\#,P_0^\#,f_0)$。还要求有另外两个等变膜 $L_k^\#=(L_k,\alpha_k^\#,P_k^\#,f_k)$，$k=1,2$，以及一个通过不变正合 Lagrangian 子流形的同痕 $\tau_V(L_1)\simeq L_2$，并可使之与等变膜结构相容。此时，可以考虑由 $L_k^\#$ 给出的 $\eps_{\A}$-不动对象 $Y_k\in\Ob\A$，以及相应的等变 Yoneda 模 $\Y_k\in\Ob\modcat^{\mathrm{eq}}(\A)$。正如\secref{6c}所指出的，扭曲 $\Ai$ 模 $T_{Y_0}(\Y_1)$ 仍然等变。\stref{定理}{17.16}的等变版本说明，$(c,k)$ 给出等变模态射
\begin{equation}\label{eq:17.24}
[t^{\mathrm{eq}}]\in\Hom_{H^0(\modcat^{\mathrm{eq}}(\A))}(T_{Y_0}(\Y_1),\Y_2),
\end{equation}
并且它在该范畴中为同构。虽然这里没有新的内容，稍作解释仍有帮助。首先，$t^{\mathrm{eq}}$ 由与\stref{引理}{5.3}(i)相同的公式定义。再回顾，若忘掉 $\Z/2$ 作用，$\A$ 就与 $\F(M)$ 拟等价。更精确地说，将每个 Lagrangian 膜赋予预等变结构，就可将 $\F(M)$ 等同于 $\A$ 的一个全 $\Ai$ 子范畴，该子范畴与整个 $\A$ 拟等价；赋予这种结构总是可能的，见\stref{笔记}{14.1}。由此得到的函子
\begin{equation}\label{eq:17.25}
\modcat^{\mathrm{eq}}(\A)\longrightarrow\modcat(\A)\longrightarrow\modcat(\F(M))
\end{equation}
在上同调意义下全忠实；第一部分忘掉等变性，在链层次上全忠实，而第二部分是向 $\F(M)\subset\A$ 的限制，由\stref{引理}{2.8}为拟等价。此外，可以适当安排嵌入 $\F(M)\subset\F(M)^{\mathrm{naive}}$，使 Lagrangian 膜 $L_k^\#$ 映到对应的等变膜。此时，$t^{\mathrm{eq}}$ 在\eref{17.25}下的像，就是前面考虑过的模同态 $t$，因此如上所述，$[t^{\mathrm{eq}}]$ 是同构。

\StatementNumber{23}\begin{lemma}\label{st:17.23}
$[t^{\mathrm{eq}}]$ 属于\eref{17.24}中 $\Hom$ 群的 $\Z/2$-不变部分。
\end{lemma}
\znote{264}{“naive”的代价正在这里：$\eps_{\A}$ 只在链同伦意义下作用，$c$ 与 $\eps_{\A}(c)$ 并不相等，仅相差一个上边界。证明正是具体给出 $(\gamma,\kappa)$，使 $\mu^1_{\A}(\gamma)=c-\eps_{\A}(c)$，并让 $k$ 与 $\eps_{\A}(k)$ 相差一个同伦，于是 $t^{\mathrm{eq}}$ 与其 $\Z/2$-像同伦。故“不变”只能在上同调层次断言，不能写成链层次的等式。[来源：Seidel, Fukaya categories and Picard–Lefschetz theory, §14e 与 §17k。]}
\begin{proof}
完全有理由认为这显然成立：此前已证明 $[t]$ 是典范的，因此它应在几何设置固有的任何对称下保持不变，等变版本也应如此。我们仍要查看细节，确保其中没有意外。可以使 Lefschetz 纤维化 $E^c\to S^c$ 与 $E^{k,r}\to S^{k,r}$，以及它们的边界条件，都带有与 $M$ 上作用相容的逐纤维 $\Z/2$ 作用。所有涉及的 Lagrangian 子流形都是等变膜，按 $\F(M)^{\mathrm{naive}}$ 的构造，它们带有 $\Z/2$-不变 Floer 数据。剩下的选择是相对扰动数据。此处，我们希望避开等变横截性问题，改作如下论证。先更一般地假设 $L_1^\#,L_2^\#$ 为预等变膜。首先考虑带有逐纤维
\SourcePage{273}
对合的 $E^c\to S^c$。以 $F$ 表示用来定义 $c$ 的 Lagrangian 边界条件，以 $F^\#$ 表示其相对膜结构，以 $(K,J)$ 表示相应的相对扰动数据。可将逐纤维对合同时作用于 $(K,J)$ 与 $F^\#$，再将后者的 Pin 结构与线丛 $\lambda^{\mathrm{top}}(TF^v)$ 作张量积。在这个改变后的情形中计数非齐次伪全纯截面，所得为
\[
CF^0\bigl(\iota_M(L_1^\#\otimes\lambda^{\mathrm{top}}(TL_1)),\iota_M(L_2^\#\otimes\lambda^{\mathrm{top}}(TL_2))\bigr)
\]
中的一个元素，按定义它就是 $\eps_{\A}(c)$。现在，若 $L_1^\#$ 与 $L_2^\#$ 均为等变膜，同构 $f_k$ 就给出自然等同
\[
CF^0\bigl(\iota_M(L_1^\#\otimes\lambda^{\mathrm{top}}(TL_1)),\iota_M(L_2^\#\otimes\lambda^{\mathrm{top}}(TL_2))\bigr)=CF^0(L_1^\#,L_2^\#),
\]
它来自 Pin 结构的同构 $P_k^\#\to\iota_M^*(P_k^\#\otimes\lambda^{\mathrm{top}}(TL_k))$。在当前情形中，这些同构与同痕 $\tau_V(L_1)\simeq L_2$ 相容，也就是说，它们给出 Lagrangian 边界条件 $F$ 的相对 Pin 结构之间的一个同构。由此可清楚看出，$\eps_{\A}(c)$ 与 $c$ 的区别仅在于相对扰动数据的选择。$k$ 的情形完全类似。因此，可以应用\eref{17.14}中的构造，插入适当符号，得到 $(\gamma,\kappa)$，满足
\[
\begin{aligned}
&\mu^1_{\A}(\gamma)+c-\eps_{\A}(c)=0,\\
&(-1)^{|c_0|-1}\bigl(\mu^1_{\A}(\kappa(c_0))-\kappa(\mu^1_{\A}(c_0))+\mu^2_{\A}(\gamma,c_0)\bigr)\\
&\qquad-\eps_{\A}(\kappa(\eps_{\A}(c_0)))+\kappa(c_0)=0
\qquad\text{对所有 }c_0\in CF^*(L_0^\#,L_1^\#).
\end{aligned}
\]
有了这些数据，\stref{引理}{5.3}(ii)的论证就给出 $t^{\mathrm{eq}}$ 与它在 $\Z/2$ 作用下的像之间的同伦，从而证明，在上同调层次上，$[t^{\mathrm{eq}}]$ 确实不变。
\end{proof}
